% Part: normal-modal-logic
% Chapter: completeness
% Section: introduction

\documentclass[../../../include/open-logic-section]{subfiles}

\begin{document}

\olfileid{nml}{com}{int}

\olsection{పరిచయం}

$\Sigma$ ఒక మోడల్ వ్యవస్థ అయితే, $\Sigma \Proves !A$ అయ్యే
ప్రతి సందర్భంలో,~$\Sigma$లోని అన్ని
\tetoken{సూత్రాల}{!!{formula}s} అన్ని నిదర్శనాలు చెల్లుబాటయ్యే
ఏ నమూనాల వర్గం~$\mClass{C}$లోనైనా $!A$ చెల్లుబాటవుతుందని
నిర్దుష్టతా సిద్ధాంతం స్థాపిస్తుంది. ముఖ్యంగా,
$\Log{K} \Proves !A$ అయితే $!A$ అన్ని నమూనాల్లో సత్యం;
$\Log{KT} \Proves !A$ అయితే $!A$ అన్ని స్వావర్తన నమూనాల్లో సత్యం;
$\Log{KD} \Proves !A$ అయితే $!A$ అన్ని సీరియల్ నమూనాల్లో సత్యం;
ఇలాగే కొనసాగుతుంది.

సంపూర్ణత నిర్దుష్టతకు విపర్యయం: ఉదాహరణకు, \Log{K}
సంపూర్ణమంటే \tetoken{ఒక సూత్రం}{!!a{formula}}~$!A$
చెల్లుబాటైతే $\Proves !A$ కావడం. సంపూర్ణతను నిరూపించడం
నిర్దుష్టతను నిరూపించడం కంటే చాలా కష్టం. ముందుగా విపర్యయ
ప్రతిజ్ఞను పరిశీలించడం ఉపయోగకరం: $\Proves/ !A$ అయినప్పుడల్లా
ఒక ప్రతినమూనా, అంటే $\mSat/{M}{!A}$ అయ్యే నమూనా~$\mModel{M}$,
ఉంటే, అప్పుడు మాత్రమే \Log{K} సంపూర్ణం. దీనికి తుల్యంగా
($!A$ను నిషేధించి), $\Proves/ \lnot !A$ అయినప్పుడల్లా
$!A$కు నమూనా ఉందని నిరూపించవచ్చు. అలాంటి నమూనా నిర్మాణంలో
$!A$లోని సమాచారాన్ని వాడుకోవచ్చు. నిర్దిష్ట
\tetoken{సూత్రాల}{!!{formula}s}కు నమూనాలు వెతికేటప్పుడూ అలాగే
చేస్తాం: ఉదాహరణకు, $p \lif \Box q$కు ప్రతినమూనా కావాలంటే,
$p$ సత్యమై $\Box q$ అసత్యమయ్యే లోకం అందులో ఉండాలి. ఒక లోకంలో
$\Box q$ అసత్యమంటే, ఆ లోకం నుంచి ప్రాప్యమైన మరో లోకంలో $q$
అసత్యమవాలి. మనకు కావలసిన సమాచారం అంతే: ఏ లోకాల్లో
\tetoken{ప్రతిజ్ఞావాక్య చరాలు}{!!{propositional variable}s}
సత్యమవుతాయి, ఏ లోకాల నుంచి ఏ లోకాలు ప్రాప్యమవుతాయి.

అయితే సంపూర్ణతను నిరూపించేటప్పుడు, నమూనా నిర్మించడానికి ఒక
నిర్దిష్ట \tetoken{సూత్రం}{!!{formula}}~$!A$ మన ముందుండదు.
$\Proves/[\Sigma] \lnot !A$ అయ్యే ప్రతి $!A$కు నమూనా ఉందని
స్థాపించాలనుకుంటాం. ఇదే కనీస అవసరం: \emph{ఒకవేళ}
$\Proves[\Sigma] \lnot !A$ అయితే, నిర్దుష్టత ప్రకారం,
$\Sigma$ సత్యమైన ఏ నమూనాలోనూ $!A$ సత్యం కాదు.
ఇప్పుడు $\Proves/[\Sigma] \lnot !A$ అయితే, అప్పుడు మాత్రమే
$!A$ $\Sigma$-అవిరుద్ధమని గమనించండి. ($\Sigma
\Proves/[\Sigma] \lnot !A$ మరియు $!A \Proves/[\Sigma] \lfalse$
తుల్యాలని గుర్తుంచుకోండి.) కనుక ప్రతి $\Sigma$-అవిరుద్ధ
\tetoken{సూత్రానికి}{!!{formula}} ఒక నమూనాను నిర్మించడం మన పని.

మన ఉపాయం $!A$ను కలిగి ఉండే $\Sigma$-అవిరుద్ధ
\tetoken{సూత్రాల}{!!{formula}s} సమితిని కనుగొనడం; అంతేకాక $!A$
సత్యమయ్యే లోకం ఎలా ఉండాలో తెలిపే ఇతర సూత్రాలనూ ఆ సమితిలో
ఉంచడం. అలాంటి సమితులే \emph{సంపూర్ణ} $\Sigma$-అవిరుద్ధ
సమితులు. $!A$ను సత్యం చేయడానికి ఒక్క లోకం గల నమూనా సరిపోదు;
దానికి అనేక లోకాలు, వాటి మధ్య ప్రాప్యత సంబంధం ఉండాలి. $!A$ను
కలిగి ఉన్న సంపూర్ణ $\Sigma$-అవిరుద్ధ సమితిలో $\Box !B$,
$\Diamond !C$ రూపాల ఇతర \tetoken{సూత్రాలు}{!!{formula}s} కూడా
ఉంటాయి. ప్రాప్యమైన అన్ని లోకాల్లో $!B$ సత్యం కావాలి;
కనీసం ఒక లోకంలో $!C$ సత్యం కావాలి. ఇది సాధించడానికి,
సాధ్యమైన \emph{అన్ని} సంపూర్ణ $\Sigma$-అవిరుద్ధ సమితులనే
లోకాల సమితికి ఆధారంగా తీసుకుంటాం. మన నమూనాలో ఒక సంపూర్ణ
$\Sigma$-అవిరుద్ధ సమితి నుంచి మరొకటి ఎప్పుడు ప్రాప్యమవుతుందో
నిర్ణయించడం కొంత కష్టమైన భాగం.

ఇలా నిర్వచించిన నమూనాలో, ఒక లోకం---అదీ ఒక సంపూర్ణ
$\Sigma$-అవిరుద్ధ సమితే---వద్ద $!A$ సత్యం కావడానికి,
$!A$ ఆ సమితిలో \tetoken{ఒక మూలకం}{!!a{element}} కావడమే
అవసరమూ చాలునూ అని చూపుతాం. $!A$ $\Sigma$-అవిరుద్ధమైతే,
అది కనీసం ఒక సంపూర్ణ $\Sigma$-అవిరుద్ధ సమితిలో
\tetoken{ఒక మూలకం}{!!a{element}} అవుతుంది (దీన్ని నిరూపిస్తాం);
కనుక $!A$ సత్యమయ్యే లోకం ఉంటుంది. అందువల్ల ప్రతి
$\Sigma$-అవిరుద్ధ \tetoken{సూత్రం}{!!{formula}}~$!A$
ఏదో ఒక లోకంలో సత్యమయ్యే ఒకే నమూనా మనకు లభిస్తుంది.
ఆ ఒకే నమూనా~$\Sigma$కు \emph{కానానికల్} నమూనా.

\end{document}
